% Part: sets-functions-relations
% Chapter: induction
% Section: induction

\documentclass[../../../include/open-logic-section]{subfiles}

\begin{document}

\olfileid[fr]{sfr}{ind}{int}

\olsection{Introduction}

\begin{explain}
La récurrence est une technique de démonstration courante en
mathématiques et en logique. Vous connaissez peut-être la récurrence
sur les entiers naturels. Mais bien d'autres ensembles d'objets
mathématiques ou logiques se prêtent à ce type de démonstration.

En général, une preuve par récurrence établit un énoncé universel :
elle montre que tous les objets d'une certaine sorte possèdent une
propriété donnée. Elle est particulièrement utile lorsque la
démonstration directe d'un énoncé universel serait trop compliquée.
Elle s'applique aux objets engendrés à partir d'objets de base par des
règles de construction : chaque objet doit pouvoir s'obtenir en un
nombre fini d'étapes de construction. Précisons cette idée.
\end{explain}

\begin{defn}
Un ensemble $S$ est \emph{stable par} une fonction $f$, définie sur
les éléments de $S$, si et seulement si, pour tout $x\in S$, on a
$f(x)\in S$.
\end{defn}

\begin{explain}
Par exemple, l'ensemble des entiers naturels ($\mathbb{N}$) est stable
par la fonction successeur $f(x)=x+1$ : si $x$ est un entier naturel,
$x+1$ en est également un. En revanche, $\mathbb{N}$ n'est pas stable
par la fonction prédécesseur $f(x)=x-1$, considérée sur les entiers
relatifs : $0$ est naturel, mais $-1$ ne l'est pas.
\end{explain}

\begin{defn}
Une propriété $P$ est dite \emph{préservée par} une fonction $f$ si,
pour tout objet $a$ du domaine de $f$, $P(a)$ entraîne $P(f(a))$ :
si $a$ possède $P$, alors $f(a)$ la possède aussi.
\end{defn}

\begin{explain}
La propriété « être un nombre pair » est préservée par la fonction
$f(x)=x+2$ : si $x$ est pair, $x+2$ l'est aussi. Elle n'est pas
préservée par la fonction successeur, puisque $x+1$ est impair quand
$x$ est pair.

La récurrence s'applique aux entiers naturels parce que chacun d'eux
s'obtient à partir de~$0$ par un nombre fini d'applications de la
fonction successeur. On procède de façon analogue pour tout ensemble
engendré à partir d'objets de base par des opérations de construction.
\end{explain}

\begin{defn}[Récurrence sur $\Nat$]
Si $0$ possède une propriété $P$ et si $P$ est préservée par la
fonction successeur, alors tout entier naturel possède $P$.
\end{defn}

\begin{ex}
La somme des entiers naturels de $0$ à $n$ est
$\frac{n(n+1)}{2}$.\\

\textbf{Cas de base.} Pour $n=0$, la somme se réduit à $0$, et
$0=\frac{0\cdot 1}{2}$.\\

\textbf{Étape de récurrence.} Supposons la propriété vraie pour $k$
et montrons-la pour $k+1$. L'hypothèse $P(k)$ s'écrit
\[ 0 + 1 + \ldots + k = \frac{k(k+1)}{2}. \]
Il faut en déduire $P(k+1)$, à savoir
\[ 0 + 1 + \ldots + k + (k+1) = \frac{(k+1)(k+2)}{2}. \]
En utilisant l'hypothèse de récurrence, on obtient
\begin{align*}
0 + 1 + \ldots + k &= \frac{k(k+1)}{2} \tag{Hypothèse} \\
0 + 1 + \ldots + k + (k+1) &= \frac{k(k+1)}{2} + (k+1)
\tag{Ajout de $k+1$ aux deux membres} \\
&= \frac{k(k+1) + 2(k+1)}{2} \\
&= \frac{k^2+3k+2}{2} \\
&= \frac{(k+1)(k+2)}{2}.
\end{align*}
L'étape de récurrence est donc établie, ce qui démontre l'énoncé.\footnote{Note
de l'édition française : une ligne intermédiaire du calcul dans la
source développe incorrectement le numérateur. L'égalité correcte
est $k(k+1)+2(k+1)=k^2+3k+2$.}
\end{ex}

\begin{explain}
Pour les !!{formula}s, les objets de base sont les !!{formula}s
atomiques, auxquels s'ajoutent les éventuelles constantes
propositionnelles primitives. Les !!{formula}s plus complexes sont
construites à partir de !!{formula}s moins complexes au moyen
d'opérateurs. Chaque opérateur correspond ainsi à une opération de
construction : la négation prend une !!{formula}, un connecteur
binaire en prend deux, et un quantificateur construit une !!{formula}
à partir d'une !!{formula} et d'une variable. Par exemple, l'opération
qui réunit $!A$ et $!B$ dans une conjonction s'écrit
$\mathcal E_\land(!A,!B)=!A\land !B$. Appelons ces opérations
des \emph{opérations de formation des formules}. L'ensemble des
!!{formula}s est stable par elles : si $!A$ et $!B$ sont des
!!{formula}s, alors $\lnot !A$, $!A\land !B$, etc., sont aussi
des !!{formula}s.
\end{explain}

\begin{defn}[Récurrence sur les !!{formula}s]
Si toutes les !!{formula}s atomiques, ainsi que les éventuelles
constantes primitives, possèdent la propriété $P$, et si $P$ est
préservée par les opérations de formation des formules, alors toute
!!{formula} possède $P$.
\end{defn}

\begin{explain}
Ce principe donne une méthode pour prouver par récurrence
qu'une propriété $P$ vaut pour toutes les !!{formula}s. On traite les
cas suivants, selon les constructeurs primitifs du langage considéré.\\

\textbf{Cas de base.} Soit $!A$ une !!{formula} atomique.
[\ldots] Donc $!A$ possède $P$. On vérifie de même les éventuelles
constantes primitives.

\textbf{Étape de récurrence.} Soient $!A$ et $!B$ des !!{formula}s
ayant la propriété $P$. Pour les opérations unaires, seule
l'hypothèse portant sur $!A$ est nécessaire.

Cas $\lnot$ : [\ldots] Donc $\lnot !A$ possède $P$.

Cas $\land$ : [\ldots] Donc $!A\land !B$ possède $P$.

Cas $\lor$ : [\ldots] Donc $!A\lor !B$ possède $P$.

Cas $\lif$ : [\ldots] Donc $!A\lif !B$ possède $P$.

Cas $\liff$ : [\ldots] Donc $!A\liff !B$ possède $P$.

Cas $\lforall[]$ : [\ldots] Donc $\lforall[x] !A$ possède $P$.

Cas $\lexists[]$ : [\ldots] Donc $\lexists[x] !A$ possède $P$.\\

Les cas correspondant aux opérations primitives étant vérifiés,
toute !!{formula} possède $P$.
\end{explain}

\end{document}
